%--------------------------------------------------------------
% abstract.tex (use thesis.tex as main file)
%--------------------------------------------------------------
\chapter*{Abstract}
%--------------------------------------------------------------

Money laundering is estimated to account for between 2\% and 5\% of global GDP annually, yet less than 1\% of illicit financial flows are successfully detected. As criminal organizations increasingly rely on complex, multi-layered transaction networks to obfuscate the origin of illicit funds, traditional rule-based and tabular machine learning approaches prove insufficient: they treat each transaction in isolation, ignoring the rich relational structure that characterizes fraudulent behavior.

This dissertation investigates the use of Graph Neural Networks (GNNs) for the automated detection of money laundering in financial transaction graphs. The problem is formulated as an edge-level classification task on the IBM AMLSim HI-Small synthetic dataset \cite{AMLWorld}, and is addressed under two different settings: a binary task distinguishing licit from illicit transactions, and a multiclass task that additionally identifies the specific laundering pattern involved, enhancing the interpretability of the model's predictions.

We propose a model based on the Group-Aware Graph Neural Network (GAGNN) architecture \cite{mainpaper}, which augments standard GNN backbones with an eMRF smoothing layer and a group representation layer that explicitly clusters accounts suspected of committing money laundering toghether. A multi-task loss combining node-level, edge-level, and group-level objectives is used to jointly optimize detection at all three levels. Two backbone variants are evaluated: one based on Graph Attention Networks (GAGNN-GAT) and one based on Graph Isomorphism Networks (GAGNN-GIN). Both are compared against the corresponding standard GNN baselines without the group-aware components.

The goal of the dissertation is to evaluate GIN as backbone for the GAGNN architecture and to understand whether the GAGNN architecture brings improvements in terms of precision, recall, F1-Score and ROC-AUC compared to baseline GNN models, either on the binary and on the multiclass task.

